% Copyright (c) 2026 Geovana Neves. Licensed under CC BY 4.0 (https://creativecommons.org/licenses/by/4.0/). See LICENSE-AND-AUTHORSHIP.md.
%
% workspaces/main.tex
%
% Guide 2 of the cfdbender guides: the workspace tree, its input files, and
% every stage of a campaign (plan, run, follow, collect and post), on a
% workstation and on an HPC cluster, then keeping and cleaning. Built by
% ../build_guides.py with this folder and ../shared/ as include directories.
%
% No em dashes and no en dashes. No underscores in frame or section titles
% outside \code: beamer writes titles to .nav and .toc. A brace group right
% after a frame title is read as the subtitle and vanishes: a frame body that
% starts with a group starts with \relax.

\documentclass[aspectratio=169,11pt]{beamer}

\input{preamble-slides.tex}
\input{hooks-slides.tex}
\cfdbinput{boxes}
\cfdbinput{hooks-contract}

\usepackage[nohyphen]{underscore}

\input{info.tex}
\guidepdftitle{2}{Workspaces}

\newcolumntype{P}[1]{>{\RaggedRight\arraybackslash}p{#1}}
\newcolumntype{R}{>{\RaggedRight\arraybackslash}X}

\begin{document}

\guidetitleframe{2}{Workspaces}{Workspaces}
  {One tree, every stage of a campaign}
  {Sources: the package pages ``Workspace layout'', ``Matrix columns'',
   ``Command line'' and ``Workspace files'' [1, 2, 3, 4].}

\begin{frame}{Who reads what}
\begin{columns}[T]
\begin{column}{0.48\textwidth}
\begin{block}{This guide}
\begin{itemize}\small
  \item Part 1: the tree, and who writes what.
  \item Part 2: the inputs you write.
  \item Parts 3 and 4: plan, then run.
  \item Parts 5 and 6: follow a run, collect its results.
  \item Part 7: a workstation and a cluster.
  \item Part 8: keeping, correcting and cleaning.
\end{itemize}
\end{block}
\end{column}
\begin{column}{0.48\textwidth}
\begin{block}{The other guides}
\begin{itemize}\small
  \item 1: the overview. Read it first.
  \item 3: solvers. The keys of each setup file, the mesh each solver reads
        and what its results mean.
  \item 4: the cheatsheet. Every command and option on one page, then one
        page per stage.
\end{itemize}
\end{block}
This guide names no solver except in a table: everything in it holds for all
six.
\end{column}
\end{columns}
\end{frame}

\begin{frame}{Contents}
\begin{multicols}{2}
\tableofcontents
\end{multicols}
\end{frame}

% Copyright (c) 2026 Geovana Neves. Licensed under CC BY 4.0 (https://creativecommons.org/licenses/by/4.0/). See LICENSE-AND-AUTHORSHIP.md.
%
% workspaces/text/01-tree-and-inputs.tex
%
% Parts 1 and 2 of guide 2: the tree and the inputs.

\section{The tree}

\begin{frame}[fragile]{One folder holds the campaign}
\begin{columns}[T]
\begin{column}{0.56\textwidth}
\begin{lstlisting}
campaign/
  init.toml              the solver of this workspace
  inputs/                yours: everything a run reads
    matrices/<name>.fs
    setups/s<id>.toml
    references/r<id>.toml
    pproc/p<id>.toml
    mesh/
  sims/                  one folder per simulation
    sim_<POL>/datapoints/DP-<point name>/
  post/<matrix>/         plan receipt, tables
  archive/<stamp>/       earlier evidence
  runs.json              the record of every point
\end{lstlisting}
\end{column}
\begin{column}{0.4\textwidth}
\begin{itemize}\small
  \item The tree is the same for every solver.
  \item A study moves to another machine by copying or syncing the folder.
  \item Nothing lives in a database or in a hidden folder of your account.
\end{itemize}
\begin{takeaway}If it is not in the folder, it is not part of the
campaign.\end{takeaway}
\end{column}
\end{columns}
\end{frame}

\begin{frame}{Who writes what}
\relax{\cfdbbodysize
\begin{tabularx}{\linewidth}{@{}P{0.27\linewidth}P{0.27\linewidth}R@{}}
\toprule
\textbf{Where} & \textbf{Written by} & \textbf{Yours to edit?} \\
\midrule
\code{init.toml} & \code{cfd-workspace init} & no \\
\code{inputs/} & you & yes, and only here \\
\code{sims/} & \code{cfd-matrix run} & no \\
\code{post/<matrix>/} & \code{plan}, \code{collect}, \code{post} & no:
  rebuilt each time \\
\code{runs.json} & \code{run}, \code{collect} & no \\
\code{marks.json} & \code{mark-failed}, \code{mark-converged} & no \\
\code{archive/<stamp>/} & any command that replaces something & no \\
\code{logs/} & \code{run}, \code{collect}, \code{post} & no \\
\bottomrule
\end{tabularx}}
\vspace{2mm}
\begin{takeaway}The files under \code{post/} are derived: deleting them loses
nothing, since \code{collect} rebuilds them from the records.\end{takeaway}
\end{frame}

\begin{frame}[fragile]{Creating a workspace}
\begin{lstlisting}
cfd-workspace init campaign --solver openfoam
\end{lstlisting}
\begin{itemize}
  \item Creates the folder if needed, \code{init.toml}, the input folders,
        \code{sims/}, \code{post/} and \code{archive/}.
  \item \code{--solver} is required. The names are lower case: \code{su2},
        \code{openfoam}, \code{pyfr}, \code{xfoil}, \code{qprop},
        \code{mock}, \code{flightstream}.
  \item A folder that already is a workspace is refused: a workspace keeps
        its solver.
  \item It also writes \code{inputs/.gitignore}, so the file that holds this
        machine's solver paths stays out of a repository.
\end{itemize}
\begin{cfdbpitfall}[One solver per workspace]
To compare two solvers on the same matrix, make two workspaces and copy the
matrix. The columns mean the same in both.
\end{cfdbpitfall}
\end{frame}

\section{The inputs}

\begin{frame}{The matrix columns}
\relax{\cfdbtablesize
\begin{tabularx}{\linewidth}{@{}P{0.25\linewidth}R@{}}
\toprule
\textbf{Column} & \textbf{Meaning} \\
\midrule
\code{POL} (required) & the simulation; its folder is \code{sims/sim_<POL>/}.
  Unique per matrix \\
\code{RUN} (required) & \code{1} plans and runs the row, \code{0} keeps it
  without running it \\
\code{SET}, \code{REF}, \code{PPROC} (required) & the setup, reference and
  post file of the row \\
\code{FLIGHT_CONDITION}, \code{SWEEP_VALUES} & the flow state and attitude;
  the values of the one key written \code{sweep} \\
\code{GEOMETRY} & the mesh or geometry file in \code{inputs/mesh/} \\
\code{AIRCRAFT}, \code{CONFIGURATION}, \code{DESCRIPTION} & your labels; they
  reach the plan listing and the summary \\
\code{NCPUS}, \code{WALLTIME} & what a cluster job asks for; \code{WALLTIME}
  carries its unit: \code{90s}, \code{240m}, \code{4h}, \code{1d} \\
\code{SOLVER_BUILD} & the solver version the row runs \\
\code{VAR_NAMES_VALUES} & per-row overrides of setup keys \\
\code{HIDDEN}, \code{SYMMETRY}, \code{SYMMETRY_LOADS}, \code{WORKFLOW} &
  kept so one layout serves every solver; the library solvers carry them or
  accept \code{-} only \\
\bottomrule
\end{tabularx}}
\vspace{1mm}
{\small An optional column may be left out of the header. A column the layout
does not know is refused before anything is planned.\par}
\end{frame}

\begin{frame}{The flight condition: the air}
\begin{columns}[T]
\begin{column}{0.5\textwidth}
\relax{\cfdbtablesize
\begin{tabularx}{\linewidth}{@{}P{0.24\linewidth}P{0.2\linewidth}R@{}}
\toprule
\textbf{Key} & \textbf{Unit} & \textbf{Role} \\
\midrule
\code{MACH} & - & velocity \\
\code{TASmps} & m/s & velocity \\
\code{REmi} & millions & density, from the Reynolds number \\
\code{ALTFT} & ft & pressure and temperature \\
\code{dISA} & K & temperature offset \\
\code{RHOkgm3} & kg/m3 & pins the density \\
\code{MUPas} & Pa s & pins the viscosity \\
\code{ASMPS} & m/s & pins the speed of sound \\
\code{TK} & K & pins the temperature \\
\code{PPA} & Pa & pins the pressure \\
\bottomrule
\end{tabularx}}
\end{column}
\begin{column}{0.46\textwidth}
\begin{itemize}\small
  \item A key carries its unit in its name.
  \item State exactly one of \code{MACH} or \code{TASmps}.
  \item No altitude means sea level of the standard atmosphere [5]; no
        \code{dISA} means zero.
  \item The viscosity follows Sutherland's law [6] unless pinned.
  \item \code{REmi} is measured against the \code{length} of the reference
        file.
  \item \code{REmi} beside \code{RHOkgm3} is refused: both fix the density.
\end{itemize}
\end{column}
\end{columns}
\end{frame}

\begin{frame}[fragile]{The flight condition: attitude and sweep}
\begin{lstlisting}
FLIGHT_CONDITION                              | SWEEP_VALUES
MACH:0.85, REmi:5, ALPHA:sweep, BETA:0        | 0,1,2,2.5,3
TASmps:40, ALTFT:5000, ALPHA:4                | -
\end{lstlisting}
\begin{itemize}\small
  \item \code{ALPHA} and \code{BETA} are in degrees. Keys match without
        regard to case; a key written twice is refused.
  \item Exactly one key may say \code{sweep}. Each value of
        \code{SWEEP_VALUES} is one \textbf{point} of the simulation.
  \item The point's folder name is built from its flight condition, so the
        same operating point has the same name in every workspace.
  \item A solver may read keys of its own, such as the rotation rate of a
        propeller. Guide 3 lists them.
\end{itemize}
\begin{cfdbpitfall}[A bare Mach number]
For some solvers a row that states only \code{MACH} and angles keeps the flow
state of its setup file. State the air in the row when you want the row to
decide it.
\end{cfdbpitfall}
\end{frame}

\begin{frame}[fragile]{Overriding a setup key in one row}
\begin{lstlisting}
VAR_NAMES_VALUES
iterations: 4000 / solver.threads: 4
\end{lstlisting}
\begin{itemize}
  \item \code{KEY: value} pairs separated by \code{ / }.
  \item Each key replaces the key of the same name in the row's setup file.
        A dotted key reaches inside a table: \code{solver.threads} is
        \code{threads} in \code{[solver]}.
  \item Use the name the setup file uses.
  \item A key that is a matrix column, such as \code{NCPUS}, is refused
        here: state it in its column.
\end{itemize}
\begin{takeaway}One setup file and a column of overrides is easier to review
than ten setup files that differ in one line.\end{takeaway}
\end{frame}

\begin{frame}{The mesh is stored once}
\begin{itemize}
  \item The \code{GEOMETRY} cell names a file in \code{inputs/mesh/}: a file
        name, or a stem that names exactly one file there.
  \item A mesh is \textbf{never copied} into a point folder. Each point
        folder holds a link to the file in \code{inputs/mesh/}.
  \item The link is relative, so it still resolves after the whole workspace
        moves to another machine.
  \item The mesh's hash is recorded with the plan and checked before the
        run.
\end{itemize}
\begin{cfdbpitfall}[Do not edit a mesh in place]
Points already run are linked to that file. A changed mesh is a new file with
a new name.
\end{cfdbpitfall}
\end{frame}

\begin{frame}[fragile]{Where the solver is: executables}
\begin{columns}[T]
\begin{column}{0.5\textwidth}
\begin{lstlisting}[style=cfdbtoml]
# inputs/executables.toml
[su2]
"8.3.0" = "/opt/su2-8.3.0/bin/SU2_CFD"
\end{lstlisting}
\begin{lstlisting}[style=cfdbtoml]
# inputs/executables.local.toml
[su2]
"8.3.0" = "D:/solvers/su2-8.3.0/SU2_CFD.exe"
\end{lstlisting}
\end{column}
\begin{column}{0.46\textwidth}
\begin{itemize}\small
  \item The first file declares the builds a row may name in
        \code{SOLVER_BUILD}. It travels with the workspace.
  \item The second gives \textbf{this machine's} paths for those builds. It
        stays out of version control.
  \item A build no file declares refuses the run before anything is
        written.
  \item Without these files, the setup's \code{[solver]} table says where
        the solver is.
\end{itemize}
\end{column}
\end{columns}
\end{frame}

% Generated inputs and saved-solution restarts.
\begin{frame}[fragile]{Generate the geometry and the mesh}
\begin{lstlisting}
cfd-geom validate model.toml
cfd-geom run model.toml
cfd-mesh airfoil-2d --naca 0012 --reynolds 6000000
\end{lstlisting}
\begin{itemize}\small
  \item \code{cfd-geom run} reaches S13 by default and records each stage under
        \code{inputs/geometry/<name>/}; \code{geometry.json} identifies it.
  \item \code{cfd-mesh mesh} meshes a configuration. \code{refine}, \code{merge}
        and the other mesh verbs transform existing files, with provenance.
  \item \code{airfoil-2d} writes \code{naca0012.su2} and \code{naca0012.msh}
        in \code{inputs/mesh/}. Name the file your solver reads in \code{MESH}.
  \item \code{--level coarse}, \code{medium} or \code{fine} chooses the cell
        counts. \code{--trailing-edge blunt} retains an airfoil base marker.
\end{itemize}
\end{frame}

\begin{frame}[fragile]{Restart uses saved state; resume selects points}
\begin{lstlisting}
cfd-matrix run polar --local --restart --sims 1001
\end{lstlisting}
\begin{itemize}\small
  \item A row's \code{RESTART} entry names a recorded source point; the
        \code{--restart} option uses each selected point's own last solution.
  \item SU2 copies its binary or ASCII restart file; OpenFOAM copies the latest
        positive time directory; PyFR copies its saved solution and native mesh.
  \item Inputs and the executable are checked before the old run is archived.
        Confirm an interrupted local process has stopped before restarting it.
        The new record keeps \code{restart_from} and the solution digests.
  \item Restart cannot combine with \code{--resume} or force-rerun options.
        Retain solution files. Other adapters warn and skip the request;
        \code{--warning-refusals} turns that skip into a refusal.
\end{itemize}
\end{frame}

% Copyright (c) 2026 Geovana Neves. Licensed under CC BY 4.0 (https://creativecommons.org/licenses/by/4.0/). See LICENSE-AND-AUTHORSHIP.md.
%
% workspaces/text/02-plan-run-follow-collect.tex
%
% Parts 3 to 6 of guide 2: plan, run, follow, collect and post.

\section{Plan}

\begin{frame}[fragile]{What plan does}
\begin{lstlisting}
cfd-matrix plan polar
\end{lstlisting}
\begin{itemize}
  \item \code{polar} is \code{inputs/matrices/polar.fs}; the \code{.fs} is
        optional.
  \item Plan reads the matrix and every file it names, resolves the air of
        each row, and asks the solver's adapter to check the inputs.
  \item It lists each point on one line and ends with the receipt and the
        next command; \code{--verbose} adds the files each run would write.
  \item It \textbf{launches nothing} and creates no point folder.
  \item It saves the plan receipt, \code{post/polar/plan.json}: the matrix
        and every input, as planned.
\end{itemize}
\begin{takeaway}Plan again after every edit of an input. \code{run} refuses a
matrix whose inputs changed since its receipt was written.\end{takeaway}
\end{frame}

\begin{frame}[fragile]{Planning part of a matrix}
\begin{lstlisting}
cfd-matrix plan polar --sims 1001 1002
cfd-matrix plan polar --sims 1001,1002
cfd-matrix plan polar --points 1001/DP-M850RE500AL+025
\end{lstlisting}
\begin{itemize}
  \item \code{--sims} takes \code{POL} values, separated by spaces or commas.
  \item \code{--points} takes point names, with or without the \code{POL/}
        prefix.
  \item A name the matrix does not carry is refused before anything
        happens.
  \item \code{run}, \code{collect} and \code{post} take the same two
        options.
\end{itemize}
\begin{block}{How much a command prints}
\code{--log-level quiet}, \code{normal} (the default) or \code{verbose};
\code{--verbose} is the short form of the last. \code{--warning-off}, written
before the command name, keeps warnings out of the console and in
\code{logs/warnings.log}.
\end{block}
\end{frame}

\section{Run}

\begin{frame}[fragile]{What run does}
\begin{lstlisting}
cfd-matrix run polar
\end{lstlisting}
\begin{itemize}
  \item Runs every active point of the planned matrix: on this machine, or
        as cluster jobs when the workspace has an HPC profile (Part 7).
  \item For each point it writes the folder
        \code{sims/sim_<POL>/datapoints/DP-<point name>/} with the solver's
        input deck, the link to the mesh, the solver's own output,
        \code{stdout.log} and \code{stderr.log}.
  \item A run on this machine also leaves \code{execution.json}, the record
        of the process that ran.
  \item It records every point in \code{runs.json}: inputs and hashes,
        status, start and finish, verdict.
\end{itemize}
\begin{cfdblimit}[The optional host guard]
\small The host guard is off by default. Add \code{--one-run-at-a-time} to
exclude other guarded local runs. Executable and memory checks always apply.
\end{cfdblimit}
\end{frame}

\begin{frame}{Running a matrix again}
\relax{\cfdbbodysize
\begin{tabularx}{\linewidth}{@{}P{0.33\linewidth}R@{}}
\toprule
\textbf{You type} & \textbf{What happens to a point already recorded} \\
\midrule
\code{run polar} & it refuses the run: finished evidence is not replaced by
  accident \\
\code{run polar --resume} & it is skipped, a failed one with the
  \code{--force-rerun} that redoes it, and the points not yet recorded run.
  A matrix that grew a row runs only the new row \\
\code{run polar --force-rerun POINT} & that point is archived, then run
  again. Repeat the option for several points \\
\code{run polar --force-rerun-all} & every recorded point of the matrix is
  archived, then run again \\
\bottomrule
\end{tabularx}}
\vspace{2mm}
\begin{itemize}\small
  \item Archived means: \code{runs.json} is copied to
        \code{archive/<stamp>/}, the point folder moves there, and the old
        record joins the history of the point. Nothing is deleted.
  \item A submitted cluster job is not canceled by a rerun: cancel it in
        the scheduler first.
\end{itemize}
\end{frame}

\section{Follow}

\begin{frame}{Six commands that only read}
\relax{\cfdbbodysize
\begin{tabularx}{\linewidth}{@{}P{0.34\linewidth}R@{}}
\toprule
\textbf{Command} & \textbf{Prints} \\
\midrule
\code{cfd-matrix status} & each simulation with the states of its points.
  \code{--points}: one row per point. \code{--failed}: only the simulations
  with a failed point \\
\code{cfd-workspace show TARGET} & one simulation, point or run record in full \\
\code{cfd-workspace history TARGET} & the current and the earlier records of a
  simulation or point, with their archives \\
\code{cfd-workspace log} & the recorded events in time order.
  \code{--problems}: the failures. \code{--since 2h}: the recent ones \\
\code{cfd-workspace trace PRODUCT} & the runs, inputs and hashes behind a
  product file \\
\code{cfd-workspace diff RUN_A RUN_B} & two run records, field by field \\
\bottomrule
\end{tabularx}}
\vspace{2mm}
\begin{itemize}\small
  \item All six take \code{--json} or \code{--csv}, to hand the rows to a
        program.
  \item They read the records in the workspace. None of them asks a
        scheduler or touches a solver.
\end{itemize}
\end{frame}

\begin{frame}{Status and verdict are two things}
\begin{columns}[T]
\begin{column}{0.48\textwidth}
\begin{block}{Status: where the point is}
\begin{itemize}\small
  \item \code{PLANNED}: registered, not started.
  \item \code{SUBMITTED}: handed to a scheduler.
  \item \code{RUNNING}: started on this machine.
  \item \code{COMPLETED}: the run ended and was read.
\end{itemize}
\end{block}
\end{column}
\begin{column}{0.48\textwidth}
\begin{block}{Verdict: how it ended}
\begin{itemize}\small
  \item \code{converged}
  \item \code{iteration_limit}
  \item \code{failed}
\end{itemize}
\end{block}
\end{column}
\end{columns}
\begin{cfdbpitfall}[Completed is not converged]
A completed point may have failed: the console shows it as
\code{FAILED_EXECUTION} or another \code{FAILED_} word, the records and tables
as \code{COMPLETED} with the verdict \code{failed}. Read the verdict, which
every table and every \code{status --points} row carries, before using a number.
\end{cfdbpitfall}
\end{frame}

\section{Collect and post}

\begin{frame}[fragile]{What collect writes}
\begin{lstlisting}
cfd-matrix collect polar
\end{lstlisting}
\relax{\cfdbtablesize
\begin{tabularx}{\linewidth}{@{}P{0.42\linewidth}R@{}}
\toprule
\textbf{File under} \code{post/polar/} & \textbf{Holds} \\
\midrule
\code{campaign_summary.json} & every recorded point with its result \\
\code{polars/FAM-ACOEF-P<POL>.csv} & coefficients per point and family,
  in wind, stability, body and rotor frames \\
\code{polars/FAM-DIM-P<POL>.csv} & dimensional loads per point and family \\
\code{polars/P<POL>-TOTAL-ACOEF.csv} & total coefficients per point \\
\code{polars/P<POL>.txt} & a 2D SU2 or OpenFOAM ALPHA sweep, XFOIL layout \\
\code{polars/P<POL>-<alias>_rotor.csv} & the rotor table, one row per point
  that turns the rotor \\
\code{polars/P<POL>_<matrix>_uns_avg.csv} & the unsteady polar: the mean over
  the averaging window of each time-accurate point \\
\code{sections/}, \code{fields/}, \code{probes/} & the products the post file
  declares, each in a folder \code{P<POL>/} \\
\code{products.json}, \code{post.log} & every product written, and every one
  skipped with the reason \\
\bottomrule
\end{tabularx}}
\begin{itemize}\small
  \item A value a point does not have is \code{NA}. No cell is quoted and
        none holds a comma, so any CSV reader takes the table.
  \item Without a matrix name, \code{collect} works on every recorded
        matrix.
  \item \code{--archive} keeps a copy of each replaced file under
        \code{archive/<stamp>/}.
\end{itemize}
\end{frame}

\begin{frame}{Products and boundary force sets}
\begin{itemize}\small
  \item Every polar has family files and family-based group files for
        \code{TOTAL}, \code{AIRFRAME} and each rotor that has members.
        Rotor solvers also write \code{FAM-RCOEF}; sectional distributions
        go under \code{sections/P<POL>/}, a volume cut under
        \code{fields/P<POL>/} and probes under \code{probes/P<POL>/}.
  \item The thrust and torque of the rotor table are counted along the
        declared shaft, so the sign of \code{CT} follows the axis. Where the
        solver supplies wall shear, the sectional loads are totals, pressure
        plus friction.
  \item Tables start with the canonical run block: conditions, air state,
        references, rotor data and averaging window. A value that does not
        apply is \code{NA}. Verdicts and histories live in the campaign summary.
  \item A setup \code{[groups]} table selects SU2 or OpenFOAM boundary markers.
        Its \code{force_groups}, including \code{TOTAL}, stay in the summary;
        they do not add boundary-group CSV files. Groups may overlap.
  \item \code{products.json}, schema version 2, states kinds, units, frames,
        families, groups, skips and SHA-256 digests. Edited or unrecorded
        products are protected against replacement. Preserve older products
        in a separate folder before recollecting.
\end{itemize}
\end{frame}

\begin{frame}{Collect or post}
\begin{columns}[T]
\begin{column}{0.48\textwidth}
\begin{block}{\code{cfd-matrix collect}}
\begin{itemize}\small
  \item Reads the saved results of each point.
  \item Picks up cluster jobs that finished.
  \item Updates the records, then rebuilds the tables.
\end{itemize}
\end{block}
\end{column}
\begin{column}{0.48\textwidth}
\begin{block}{\code{cfd-matrix post}}
\begin{itemize}\small
  \item Rebuilds the same tables from the records and the saved outputs
        alone.
  \item Runs nothing and asks nothing of a cluster.
\end{itemize}
\end{block}
\end{column}
\end{columns}
\vspace{1mm}
\relax{\cfdbbodysize
\begin{tabularx}{\linewidth}{@{}P{0.16\linewidth}R@{}}
\toprule
\textbf{Exit code} & \textbf{Meaning} \\
\midrule
0 & done \\
1 & done, and at least one point failed (\code{diff}: the two runs differ) \\
2 & refused: the message names the file and the cell or key at fault, and
  nothing was written \\
3 & done, and submitted work is still outstanding \\
\bottomrule
\end{tabularx}}
\end{frame}

% Copyright (c) 2026 Geovana Neves. Licensed under CC BY 4.0 (https://creativecommons.org/licenses/by/4.0/). See LICENSE-AND-AUTHORSHIP.md.
%
% workspaces/text/03-cluster-and-keeping.tex
%
% Parts 7 and 8 of guide 2: a workstation and a cluster; keeping, correcting
% and cleaning.

\section{A workstation and a cluster}

\begin{frame}{Two machines, two jobs}
\begin{columns}[T]
\begin{column}{0.48\textwidth}
\begin{block}{The workstation}
\begin{itemize}\small
  \item Write the inputs and plan.
  \item Run one point, or a cheap solver, and look at it.
  \item Collect, read the tables, plot.
\end{itemize}
\end{block}
\end{column}
\begin{column}{0.48\textwidth}
\begin{block}{The cluster}
\begin{itemize}\small
  \item Run the long matrix, one scheduler job per point.
  \item Keep the heavy solver output where it was written.
\end{itemize}
\end{block}
\end{column}
\end{columns}
\begin{itemize}
  \item The workspace is the same tree on both. What differs is one file,
        the \textbf{HPC profile}, that says how this site submits a job.
  \item A workspace with a profile submits; \code{--local} keeps the work on
        this machine all the same.
\end{itemize}
\end{frame}

\begin{frame}[fragile]{The HPC profile}
\begin{columns}[T]
\begin{column}{0.5\textwidth}
\begin{lstlisting}[style=cfdbtoml]
# inputs/hpc/h001.toml
application_id = "SU2"
max_walltime = "96:00:00"
[descriptor]
format = "yaml"
name = "submit.yaml"
[descriptor.fields]
job_name = "{job_name}"
walltime = "{walltime}"
ncpus = "{ncpus}"
version = "{solver_version_alias}"
working_directory = "{workdir}"
config_file = "{config_path}"
[submit]
command = ["submit-tool", "{descriptor_path}"]
[builds]
"8.3.0" = "8.3.0"
\end{lstlisting}
\end{column}
\begin{column}{0.46\textwidth}
\begin{itemize}\small
  \item The profile describes a job file and the command that submits it.
        \code{submit-tool} stands for your site's own.
  \item \code{\{ncpus\}} and \code{\{walltime\}} come from the row's
        \code{NCPUS} and \code{WALLTIME}.
  \item \code{[builds]} maps the row's \code{SOLVER_BUILD} to the name the
        site gives that version.
  \item \code{\{config_path\}} is the file the solver reads in the point
        folder.
  \item \code{--hpc NAME} chooses the default; a row's \code{hpc: NAME}
        in \code{VAR_NAMES_VALUES} overrides it.
\end{itemize}
\end{column}
\end{columns}
\end{frame}

\begin{frame}[fragile]{Select a profile and cluster per row}
\begin{lstlisting}
VAR_NAMES_VALUES
hpc: alpha / cluster: west
hpc: beta / cluster: east
\end{lstlisting}
\begin{itemize}\small
  \item \code{hpc: alpha} selects \code{inputs/hpc/alpha.toml} before the
        \code{--hpc} default. A profile and its digest are recorded per point.
  \item \code{cluster: west} supplies the profile's \code{\{cluster\}}
        placeholder; it falls back to \code{cluster} in \code{init.toml}.
        Keys are case-insensitive and never passed to the solver setup.
  \item \code{--local} ignores both. An unused cluster selection warns;
        \code{--warning-refusals} makes that a refusal.
  \item Plan again after changing a profile. Older HPC receipts without the
        selected profile and digest are refused. Collect waits for all
        configured \code{job_end_files} and settled outputs.
\end{itemize}
\end{frame}

\begin{frame}{A cluster run, step by step}
\begin{enumerate}
  \item \code{cfd-matrix plan polar} prints each point, its job file and the
        submit command. Nothing is written but the receipt.
  \item \code{cfd-matrix run polar} stages every file the solver reads into
        the point folder, then submits one job per point. The points are
        recorded \code{SUBMITTED}.
  \item The job runs in its point folder and needs nothing else from the
        workspace.
  \item \code{cfd-matrix collect polar} reads the points whose output files
        exist and leaves the others outstanding. \code{--watch} repeats the
        sweep until none is outstanding.
\end{enumerate}
\begin{cfdbpitfall}[Send the job's output to the point folder]
Point the scheduler's output files at \code{\{stdout_log\}} and
\code{\{stderr_log\}}: several solvers print their results there, and
collect reads them from the point folder.
\end{cfdbpitfall}
\end{frame}

\begin{frame}[fragile]{Bringing results home: sync}
\begin{lstlisting}
cfd-workspace sync runs
cfd-workspace sync post --apply
cfd-workspace sync all --apply
\end{lstlisting}
\begin{itemize}
  \item \code{inputs/sync-workspaces.toml} names the other workspaces;
        \code{--from NAME} takes one of them.
  \item Each level adds to the one before: \code{runs}, then \code{post},
        then \code{all}.
  \item Sync \textbf{previews} until \code{--apply}.
  \item On a conflict it keeps this workspace's record and file, unless
        \code{--prefer-other} (records) or \code{--overwrite} (files, the
        local copy archived first) says otherwise.
\end{itemize}
\begin{takeaway}Bring the records and the tables to the workstation. Leave
the heavy files on the cluster until you need one.\end{takeaway}
\end{frame}

\section{Keeping, correcting, cleaning}

\begin{frame}{Correcting the records}
\relax{\cfdbbodysize
\begin{tabularx}{\linewidth}{@{}P{0.4\linewidth}R@{}}
\toprule
\textbf{Command} & \textbf{Does} \\
\midrule
\code{cfd-workspace mark-failed --sims POL} & records your verdict that points
  failed, whatever the solver reported \\
\code{cfd-workspace mark-converged --sims POL --reason TEXT} & records your
  verdict that completed points converged; the reason is required \\
\code{cfd-workspace restore KIND} & brings \code{runs}, \code{storage},
  \code{products} or \code{plan} back from \code{archive/<stamp>/} \\
\code{cfd-workspace rebuild} & writes records for simulation folders that have
  none \\
\code{cfd-workspace archive SIM_ID} & moves one simulation's folder and
  records to \code{archive/<stamp>/} \\
\bottomrule
\end{tabularx}}
\vspace{2mm}
\begin{itemize}\small
  \item The correction commands preview until \code{--apply}.
  \item \code{cfd-workspace archive} acts at once, without a preview.
  \item Your verdicts are kept in \code{marks.json}, apart from what the
        solver reported.
\end{itemize}
\end{frame}

\begin{frame}{Space on disk}
\relax{\cfdbbodysize
\begin{tabularx}{\linewidth}{@{}P{0.4\linewidth}R@{}}
\toprule
\textbf{Command} & \textbf{Does} \\
\midrule
\code{cfd-workspace space-in-use} & sizes by folder, by simulation and by file
  extension \\
\code{cfd-workspace free-space RECIPE} & deletes the files that
  \code{inputs/management/<RECIPE>.toml} names, except those a record needs \\
\code{cfd-workspace delete-sims --sims POL} & deletes simulation folders and
  retires their records \\
\bottomrule
\end{tabularx}}
\vspace{2mm}
\begin{itemize}\small
  \item \code{free-space} and \code{delete-sims} preview until
        \code{--apply}. Read the preview: these are the only commands that
        delete.
  \item Every storage command is recorded in \code{storage_management.json};
        \code{cfd-workspace log --storage} lists them.
\end{itemize}
\end{frame}

\begin{frame}{Version control}
\begin{columns}[T]
\begin{column}{0.48\textwidth}
\begin{block}{Worth committing}
\begin{itemize}\small
  \item \code{init.toml}
  \item \code{inputs/}: small text, and the whole definition of the study
  \item \code{runs.json} and \code{post/}, when the results are part of the
        record you keep
\end{itemize}
\end{block}
\end{column}
\begin{column}{0.48\textwidth}
\begin{block}{Usually left out}
\begin{itemize}\small
  \item \code{sims/} and \code{archive/}: solver output
  \item large meshes: keep them where your site keeps large files
  \item \code{inputs/executables.local.toml}: one machine's paths, ignored
        by the \code{.gitignore} that \code{init} writes
\end{itemize}
\end{block}
\end{column}
\end{columns}
\begin{takeaway}The inputs and the records are enough to say what was run.
The solver output is what it cost to run it.\end{takeaway}
\end{frame}


\begin{frame}{References}
\begin{reflist}
  \bibitem{g2r01} \docref{workspace-layout}{Workspace layout}
  \bibitem{g2r02} \docref{matrix-columns}{Matrix columns}
  \bibitem{g2r03} \docref{cli}{Command line}
  \bibitem{g2r04} \docref{workspace-layout}{Workspace layout}
  \bibitem{g2r05} International Organization for Standardization, \emph{ISO
        2533:1975, Standard Atmosphere}, ISO, Geneva, 1975.
  \bibitem{g2r06} F. M. White, \emph{Viscous Fluid Flow}, 3rd ed.,
        McGraw-Hill, New York, 2006. Sutherland's law for the viscosity of
        air.
\end{reflist}
\end{frame}

\end{document}
